% BEGIN APPROVED GEOMETRIC FUBINI
\begin{nota}\label{nota:geometric:fubini:target}
For a parameter base $\mathbb T$, put
\[
 R(\mathbb T)=\sggk\RES_{\mathbb T},\qquad
 R_g(\mathbb T)=R(\mathbb T)[L^{-1}],\qquad L=[\A].
\]
We use $\bb E_g^{\mathbb T}$ for the geometric retraction of Section~\ref{sec:retract}
followed by the natural homomorphism to $R_g(\mathbb T)$.
The RES grading is thus forgotten in the target. Write $\widetilde{\bb E}_g^{\mathbb T}$
for the unnormalized geometric Euler homomorphism followed by forgetting
the RES grading. If $\bm U\in\RV[k]$, then
\begin{equation}\label{geometric:fubini:normalization}
 \bb E_g^{\mathbb T}([\bm U])=L^{-k}\widetilde{\bb E}_g^{\mathbb T}([\bm U]).
\end{equation}
The class $L$ need not be a non-zero-divisor in $R(\mathbb T)$.

\end{nota}

\begin{prop}[Geometric version of Theorem~\ref{fubini}]\label{prop:geometric:fubini}
Let $\mathbb S$ be a $(\VF,\Gamma)$-generated ground substructure,
with $\Gamma(\mathbb S)$ nontrivial. Let $A,A'\in\VF_*$ and let
$f:A\to B$, $f':A'\to B$ be definable maps. Put
$A_b=f^{-1}(b)$, $A'_b=(f')^{-1}(b)$ and
$\mathbb S_b=\mathbb S\langle b\rangle$.
Suppose that for every $b\in B$ there are
$\mathbb S_b$-definable objects $\bm U_b,\bm U'_b\in\RV[*]$ and
$\mathbb S_b$-definable bijections
\[
 A_b\xrightarrow{\ \sim\ }\bb L\bm U_b,
 \qquad A'_b\xrightarrow{\ \sim\ }\bb L\bm U'_b
\]
such that
\[
 \bb E_g^{\mathbb S_b}([\bm U_b])
 =\bb E_g^{\mathbb S_b}([\bm U'_b])
 \quad\text{in }R_g(\mathbb S_b).
\]
Then
\[
 \bb E_g^{\mathbb S}\!\left(\int[A]\right)
 =\bb E_g^{\mathbb S}\!\left(\int[A']\right)
 \quad\text{in }R_g(\mathbb S).
\]
As in Theorem~\ref{fubini}, the hypothesis concerns all points of $B$ in the
sufficiently saturated ambient model.
\end{prop}

\begin{proof}
We explain the changes to the proof of Theorem~\ref{fubini}, retaining its
positive lifting witnesses and intrinsic parameter-relative RES quotients.

Fix $b$ temporarily, and let $k,k'$ be the degrees of its two RV
representatives. Write
\[
 a=\widetilde{\bb E}_g^{\mathbb S_b}([\bm U_b]),\qquad
 a'=\widetilde{\bb E}_g^{\mathbb S_b}([\bm U'_b]).
\]
By the definition of localization, the equality $L^{-k}a=L^{-k'}a'$
is equivalent to the existence of an integer $N\geq k,k'$ such that
\begin{equation}\label{geometric:fubini:clearing}
 L^{N-k}a=L^{N-k'}a'\quad\text{in }R(\mathbb S_b).
\end{equation}
Indeed, after a common denominator has been chosen, some further power
of $L$ annihilates the difference of the two numerators. This step does
not cancel $L$ in the unlocalized ring.

Let $\bm{\mathbb A}_{\RES}\in\RES[1]$ be the disjoint union of
$\G_m$ and the degree-one point. Its ungraded class is $L$.
Its lift is $\mathcal O^\times\sqcup(1+\mathcal M)$, which is definably
bijective to the closed unit ball $\mathcal O$, by translating the second
summand to $\mathcal M$. Consequently
\[
 \bb E_g([\bm{\mathbb A}_{\RES}])=L^{-1}L=1.
\]
Replace the two representatives by
\[
 \widehat{\bm U}_b=\bm U_b\times\bm{\mathbb A}_{\RES}^{\,N-k},
 \qquad
 \widehat{\bm U}'_b=\bm U'_b\times\bm{\mathbb A}_{\RES}^{\,N-k'}.
\]
Both have degree $N$, and \eqref{geometric:fubini:clearing} says that their $\widetilde{\bb E}_g$-values
are equal. Their VF lifts are definably bijective to
$A_b\times\mathcal O^{N-k}$ and $A'_b\times\mathcal O^{N-k'}$.

The integer $N$ is part of the finite witness package. More precisely,
express \eqref{geometric:fubini:clearing} by the finite stabilized RES-bijection criterion
used in Theorem~\ref{fubini}, and include the padding exponents, all its objects,
bounds and graphs, and both positive lifting bijections. The validity of
each fixed package is first order over the original ground. These
conditions cover every ambient parameter $b$. Compactness supplies
finitely many packages; a finite definable partition fixes $k,k',N$ and
the witness formulas on each part. No uniform bound on the
denominator-clearing exponents is assumed before this argument.

Work on one such part. The relative RV calculation in the proof of
Theorem~\ref{fubini} applies to $\widetilde{\bb E}_g$ with $\chi_g$ in place of $\chi_b$.
For clarity, if $X$ is a definable RV set,
$J\subseteq X\times\Gamma^r$ is definable, and
\[
 Y=\{(x,z)\in X\times(\RV^\times)^r:\vrv(z)\in J_x\},
\]
then its unnormalized form is
\begin{equation}\label{geometric:fubini:push}
 \widetilde{\bb E}_g([Y])=g^r\sum_e e\,\widetilde{\bb E}_g([X_e]),\qquad
 g=[\G_m],\quad X_e=\{x:\chi_g(J_x)=e\}.
\end{equation}
This follows from the same monomial decomposition and Fubini formula
for $\chi_g$ on the $\Gamma$-sort. The sum is finite, and its level sets are definable over the
original ground. For the first-order Euler certificates, use finite
decompositions into relatively open rational polyhedra: a nonempty
cell of dimension $j$ has geometric Euler characteristic $(-1)^j$.
The cell formulas, dimensions and nonemptiness flags are included in
the finite package. The varying value-torsor identity also remains valid,
since $\chi_g$ of a singleton is $1$. Thus the actual total-bijection
relations and the value-fiber relations in the finite stabilized
witnesses are preserved by $\widetilde{\bb E}_g$. Their common stabilization is
cancelled additively, as in the proof of Theorem~\ref{fubini}.

Include all these formulas in the same preparation package. On a
prepared base piece, let $d$ be the number of nonzero VF base coordinates
and write $\tau=(s,t)$ as in that proof. The padded total RV objects have
structure maps
\[
 q(\tau,u)=(s,p_\tau(u)),\qquad
 q'(\tau,u')=(s,p'_\tau(u')),
\]
and both belong to $\RV[d+N]$. Finiteness of these maps follows from the
actual RV fibers of the total VF lifts, as before. In particular, the
retained coordinates $t$ receive no additional VF lifts. The relative
calculation and complete stabilized witnesses give equality of the
$\widetilde{\bb E}_g$-values of the two total objects. Multiplication by the common
factor $L^{-(d+N)}$ and \eqref{geometric:fubini:normalization} give equality of their
geometric values.

On the chosen part the padded total VF spaces are
$A|_{B_i}\times\mathcal O^{N-k}$ and $A'|_{B_i}\times\mathcal O^{N-k'}$.
Since $\bb E_g(\int[\mathcal O])=1$, their geometric values agree with those of
$A|_{B_i}$ and $A'|_{B_i}$, respectively. Finite additivity over the base
partition proves the assertion.
\end{proof}

\begin{rem}
Taking $A'=\varnothing$ gives the corresponding vanishing statement.
The proposition is stated in the ungraded localized target
$R_g(\mathbb S)$. It uses no injectivity of the natural homomorphism
from the homogeneous degree-zero affine localization to that target.
\end{rem}

{\normalfont\small\color{ResolutionGreen}
\noindent\textbf{Author resolution (geometric Fubini).}
The author approved this geometric version on November 14, 2026,
including its ungraded affine-localized target, denominator-clearing
witnesses before preparation, affine-line padding and common total-degree
normalization. It is added to the Section 5 formalization scope.\par}
% END APPROVED GEOMETRIC FUBINI

